% Ported verbatim from paper/top3_joint_moment_bound.tex lines 1280-1346.
% Only cross-reference wording was adapted; no mathematics was altered.
By Theorem~\ref{structure-main}, write the canonical law as
$\sum_{i=1}^m p_i\delta_{v_i}$ with $m\le24$ and all $p_i>0$.
Set $R_i=|v_i|$, $u_i=v_i/R_i$, $w_{ij}=u_i\cdot u_j$, and
\[
 \kappa=\frac3{125000},\qquad
 h(t)=\sigma(t)-\frac{z_*(t)}{125}+\kappa,
 \qquad \chi=\sum_i p_i^2R_i^2.
\]
Equation~\eqref{contact-paid} gives $h\ge0$ and $h(1)=\kappa$.
The same contact remainder in~\eqref{Cstar8} therefore has the exact split
\begin{equation}\label{atom-split}
 r_{\rm ct}=\kappa\chi+r_{\rm off},\qquad
 r_{\rm off}=\sum_{i\ne j}p_ip_jR_iR_jh(w_{ij})\ge0.
\end{equation}
These self-pairs have positive product mass. They are not values assigned
to the null diagonal $x=y$ of the underlying nonatomic space.
The estimate $\chi\ge C^2/m$ is the usual Cauchy--Schwarz
cardinality bound. For related weighted-energy arguments,
see~\cite[Theorem 2.1 and the subsequent discussion]{Weighted}.

In the branch $D_*<1/20$, let
\[
 \widehat Q=Q-\frac{V}{9000000}.
\]
It is nonnegative: in~\eqref{Q8} only the last separate $V$ coefficient
changes, to
\[
 \frac{42847433}{5229378000000}-\frac1{9000000}
       =\frac{14088797}{1743126000000}>0.
\]
\begin{proposition}\label{final-certificate}
For every maximizing graphon with $D_*<1/20$,
\begin{align}\label{certificate-final}
 D_*={}&d+\widehat Q+{\cal R}_0+\ell r_{\rm reg}
       +\frac{K^2}{1125}+r_{\rm tr}\notag\\
 &+\frac{\kappa}{9}\left(\chi-\frac{C^2}{24}\right)
       +\frac{r_{\rm off}}9.
\end{align}
Every term after $d$ is nonnegative. Consequently $\mathcal M\le c$.
\end{proposition}
\begin{proof}
Lemma~\ref{spectral-reserve} verifies the radial hypothesis of
Proposition~\ref{conditional8}. Substitute~\eqref{atom-split} once into
\eqref{certificate8}. Since $C^2=3-V$,
\[
 \frac{\kappa\chi}{9}
 =\frac1{3000000}-\frac{V}{9000000}
       +\frac{\kappa}{9}\left(\chi-\frac{C^2}{24}\right).
\]
This gives~\eqref{certificate-final}. Moreover
\[
 \chi-\frac{C^2}{24}
 =\sum_i\left(p_iR_i-\frac Cm\right)^2
       +C^2\left(\frac1m-\frac1{24}\right)\ge0.
\]
No lower bound on an individual mass is needed. The complementary branch
also gives $D_*\ge d$, since the rational estimates used above imply
\[
 \frac1{20}-d>
       \frac{1025023913201}{804826176000000}>0.
\]
Thus $D_*\ge d$ in both branches, or equivalently $\mathcal M\le c$.
\end{proof}
The atomic term changes the coefficient by $1/6000000$ relative to
$(3+\sqrt5)/4-d_1/2$. Most of the improvement in Theorem~\ref{main}
comes from the preceding joint moment estimate, not from this last term.
